% Contenido incluido desde la sección de justificación.
\subsection{Viabilidad}
El uso de códigos QR para el control de acceso físico ya se ha propuesto en
la literatura \parencite{kao2011physical}, y su bajo costo frente a otras
tecnologías (Tabla~\ref{tab:comparativa}) hace viable construir y probar un
prototipo con los recursos del equipo. Frente a NFC, una comparación en
transacciones móviles encontró que el QR se procesa con mayor rapidez, aunque
no incorpora cifrado propio \parencite{harnaningrum2024comparison}; esa
carencia se compensa con la firma HMAC descrita en Delimitación. Además, el
estudiante usa su propio celular, por lo que no hace falta emitir tarjetas.
Frente al reconocimiento facial, el QR no requiere procesamiento pesado ni
trata datos biométricos, cuya aceptación entre estudiantes universitarios
depende en buena medida de la confianza y de la seguridad que perciben
\parencite{rukhiran2023user}, y cuyo tratamiento implicaría, además, las
consideraciones de protección de datos personales abordadas en la metodología.

\begin{table}[tbp]
	\caption{Comparación de Tecnologías de Validación de Acceso}
	\label{tab:comparativa}
	\begin{tabular}{lccc}
		\toprule
		Criterio                            & Código QR                & NFC     & Reconocimiento facial \\
		\midrule
		Costo de infraestructura            & Bajo                     & Medio   & Alto                  \\
		Medio que porta el estudiante       & Celular                  & Tarjeta & Ninguno               \\
		Complejidad de cómputo en el lector & Baja                     & Baja    & Alta                  \\
		Trata datos biométricos             & No                       & No      & Sí                    \\
		Riesgo de copia del identificador   & Medio\textsuperscript{a} & Bajo    & Bajo                  \\
		\bottomrule
	\end{tabular}
	\tablenote{Elaboración propia con base en
		\textcite{harnaningrum2024comparison} y \textcite{rukhiran2023user}. El costo
		y la complejidad son una valoración cualitativa del equipo.
		\textsuperscript{a}El riesgo corresponde a un QR estático; con un código
		dinámico firmado con HMAC se reduce considerablemente
		\parencite{subpratatsavee2015hmac}.}
\end{table}

Queda por conocer la viabilidad tecnológica desde el lado del estudiante, es
decir, si cuenta con un celular propio, con conexión a internet y con batería
suficiente al llegar al plantel; esto se indaga en la Sección~IV del
Apéndice~A.

La investigación es necesaria porque hoy no existen datos que midan el
problema en el \tescha{}: se desconoce cuánto tiempo espera un estudiante para
entrar y con qué frecuencia se usan credenciales ajenas. Levantar esa información y
contrastarla con el desempeño de un prototipo permitirá saber, con evidencia y
no con suposiciones, si esta solución de bajo costo ayuda a aliviar los
estragos del sistema actual y si justifica su adopción en el resto de los
accesos del plantel.
